\chapter*{Before you start}
\addcontentsline{toc}{chapter}{Before you start}
\markboth{BEFORE YOU START}{BEFORE YOU START}

Chris,

this document exists because I kept the work to myself for too long. Between the August pilots and
this morning I ran a long series of measurements on our breakdown-radius line. A few pieces reached
you: the evaluation plan on the branch \texttt{iclr2027-eval-plan}, and what we said to each other on
3 September. Most of it did not. The August package waited on a branch I never pushed, and
everything from 7 to 16 September lives in folders on my machine. Your draft of 16 September lists,
among the experiments it still owes, the ranking evaluation on the real instances and the elicitation
arm. I ran the elicitation arm between 11 and 13 September, and a corrupted-knowledge calibration on
the same real rows, though not your $\tau_b$ design itself (Chapter~\ref{ch:merge} is precise about
the difference). That overlap is the clearest sign that the gap is mine to close, and this close to the
author lock is late to close it.

So here is all of it, in two forms. The code, the results and the pre-registrations are on one
branch of \texttt{bkrobust}, \texttt{ze/experiments}, built on top of your history so that it merges.
This polycopié is the explanation: what I ran, why, what came out, what I had to retract, and what it
changes in the paper.

\begin{ideia}[The whole document in four sentences]
Knowledge buys identifiability, then robustness, never efficiency, and one wrong claim usually costs
the whole of what it bought. On the true CPDAG, counted in claims, the certificate never
over-certified on 3,014 certifiable rows of eleven knowledge conditions; counted in covering hops, it
did on 39. On the one panel where the class itself was estimated, no retraction reached the truth on
any certifiable row, so the guarantee is conditional on a correct equivalence class. The
paper that survives is a sensitivity analysis for background knowledge in the sense of Rosenbaum's
$\Gamma$, whose honest baseline is a leave-$k$-out screen that it ties.
\end{ideia}

\section*{How it is organised}

\textbf{Part~I} is short. It states the object once in your geometry and gives the translation table
to the words I use everywhere else: claim, retraction, hop, $\rclaim$, $\rhop$, $\rnul$, the
knowledge interval. If you read only one chapter before our call, read that one and
Chapter~\ref{ch:merge}.

\textbf{Parts~II to~V} are the measurements, in the order they happened. August (the pilot, the
\texttt{se} criterion, locality, the firewall census, the wrong-skeleton experiment). Early September
(what the radius should measure, and three defects in the old corpus). The protocol executed on
11--13 September (the observable frame, the knowledge suppliers, Table~4, the learned CPDAG). And
the reframing as a Rosenbaum analogue, with the two measurements that decide it (M1, M1b) and the
tables of 16 September.

\textbf{Part~VI} is your draft. Chapter~\ref{ch:review} reads it against the measurements: what is
right, what the text says that the mathematics does not, and a handful of mechanical defects visible
in the PDF. Chapter~\ref{ch:merge} is the merge: the four conflicts, three options, the one I
recommend, and what has to happen by Friday.

\textbf{Part~VII} maps the repository, file by file, with how to reproduce each table.

\section*{How to read the boxes}

Green boxes (\emph{What was measured}) hold numbers computed on my machine, and each names the
repository file or script that produced it. Orange boxes point to a paper worth reading. Red boxes
are errors actually made, ours or the draft's, with what caught them. Purple boxes are sentences in
the form I think the paper can print. Aqua boxes are commands and decisions.

The pink boxes are questions to answer without looking; the answers are in the appendix. They are
there because this is a lot of material to absorb before Friday, and answering a question fixes a
number better than rereading it. Skip them if you prefer; nothing in the argument depends on them.

\digressaoaviso

\section*{What I am confident about, and what I am not}

Every number in the green boxes was produced by a script and re-read from its output file while
writing; the few whose script is not on the branch say so in the box. Several numbers I had quoted earlier turned out to be wrong, and the chapters
say which, when and how they were caught; I would rather you see the corrections than meet them in
review. The theory of your draft was checked by a re-implementation of the space, the order, the covering
relation and the distance from your definitions, which reuses your Meek-closure and extension code,
and no counterexample was found. The numbers printed in your draft were not audited.

One file in the repository is older than its data: \texttt{results/stage0/STAGE0.md} quotes the
calibration run from before a seeding fix, while the calibration files next to it are the rerun
(Chapter~\ref{ch:frame} explains). Chapter~\ref{ch:repo} lists this and every other caveat about
the files.

Zé \hfill 16 September 2026
